\section{Related Work}

\vspace{-0.5em}
\subsection{CLIP for 3D Learning}
\vspace{-0.5em}

Image-language models like CLIP have achieved remarkable performance through large-scale image-text pretraining~\cite{radford2021learning,jia2021scaling,li2022grounded,zhang2022glipv2,alayrac2022flamingo,ramesh2022hierarchical,saharia2022photorealistic}. As these models excel at capturing rich visual concepts and possess impressive zero-shot capabilities, they have been applied to various 3D vision tasks. For instance, numerous recent works utilize CLIP to facilitate zero-shot text-to-3D generation~\cite{hong2022avatarclip,jain2022zero,michel2022text2mesh,sanghi2022clip,lee2022understanding,canfes2023text,khalid2022text,aneja2022clipface,jetchev2021clipmatrix,xu2022dream3d,liu2022iss}, typically through CLIP-guided per-scene optimization. From a recognition perspective, some works focus on scene-level representation, aiming to leverage CLIP priors for zero-shot 3D segmentation or detection in both indoor~\cite{ha2022semantic,jatavallabhula2023conceptfusion,ding2022language,yang2023regionplc,lu2023open,peng2022openscene,zeng2023clip,huang2023joint,rozenberszki2022language,zhang2023clip,kerr2023lerf} and outdoor scenes~\cite{chen2023clip2scene,hess2022lidarclip}. Meanwhile, another line of work focuses on shape-level understanding, targeting zero-shot shape classification~\cite{zhang2022pointclip,zhu2022pointclipv2,qi2023contrast,xue2022ulip,hegde2023clip} and part segmentation~\cite{liu2022partslip,abdelreheem2023satr}. There are two primary working paradigms for these methods. The first~\cite{zhang2022pointclip,zhu2022pointclipv2,huang2022clip2point} involves using images as a medium representation, projecting 3D point clouds into 2D and employing 2D CLIP for inference. However, these methods typically suffer from occlusion and information loss during projection, along with unnecessary latency due to point cloud rendering and multiple 2D CLIP inferences. The second paradigm involves training a 3D-native encoder attempting to distill or fuse CLIP features into 3D representations. Our paper follows this paradigm.

\vspace{-0.5em}
\subsection{3D Shape Representation Learning}
\vspace{-0.5em}

Various works have studied self-supervised pretraining for point clouds by designing pretext tasks~\cite{eckart2021self,thabet2020self,poursaeed2020self,achituve2021self,sharma2020self} such as self-reconstruction~\cite{rao2020global,deng2018ppf,achlioptas2018learning,wang2021unsupervised}, masked auto-encoding~\cite{pang2022masked,yu2022point,hess2023masked}, distortion reconstruction~\cite{sauder2019self,mersch2022self,sun2021point}, normal estimation~\cite{rao2020global}, and contrastive learning~\cite{zhang2021self,sanghi2020info3d,xie2020pointcontrast}. These tasks enhance models' shape representations and improve their performance on downstream applications, although they do not involve multimodal semantic alignments during pretraining. 

Recently, some works~\cite{qi2023contrast,xue2022ulip,hegde2023clip}, exemplified by ULIP~\cite{xue2022ulip}, have explored learning multimodal joint representations for 3D shapes. They train 3D-native shape encoders by aligning 3D shape embeddings with CLIP's language and/or image embeddings through multimodal contrastive learning. Works like ReCon~\cite{qi2023contrast} further combines cross-modal contrastive learning with masked auto-encoding for added enhancement. While these methods allow for zero-shot 3D classification through the computation of 3D-text similarity, the amount of distilled knowledge and their model capability are heavily limited by the small-scale training datasets used. Our work follows this paradigm but aims to learn more generalizable and scalable representations to enable open-world 3D shape understanding.
